\chapter{Phase 4: Production and Safety (The Harness Phase)}

The final phase of becoming a Harness Engineer involves deploying agents safely and measuring their performance. We implement the complete safety harness: sandboxed computation, automated evaluation, and human approval gates.

\section{Step 10: Sandboxed Code Execution}
Executing agent-generated code requires an isolated environment. Anthropic's \texttt{code\_execution} tool provides a secure, server-side container with popular data science libraries pre-installed.

\begin{lstlisting}[caption={step10\_sandboxed\_agent.py}, label={lst:step10}]
"""
PHASE 4 - STEP 10: Sandboxed Code Execution
Harness Pillar: ACI (Compute Layer)
Goal: Securely execute agent-generated Python code.
"""

import anthropic

def run_sandboxed_agent(task: str):
    response = client.beta.messages.create(
        model="claude-opus-4-5-20251101",
        max_tokens=4096,
        betas=["code-execution-2025-08-25"],
        tools=[{"type": "code_execution_20250825"}],
        system="You are a data analyst. Write and execute Python code.",
        messages=[{"role": "user", "content": task}],
    )
    return response.content[0].text
\end{lstlisting}

\subsection*{Key Takeaways}
\begin{itemize}
    \item Use isolated containers with no network access for agent code execution to prevent data exfiltration.
    \item Anthropic's sandbox comes with pre-installed data science tools like \texttt{pandas} and \texttt{matplotlib}.
    \item If running your own sandbox (e.g., Docker), always use \texttt{--network none} and \texttt{--read-only} flags.
\end{itemize}

\section{Step 11: Evaluation Suite}
Automated evaluation is critical for measuring agent performance. We use the ``LLM-as-judge'' pattern to grade agent outputs against a test dataset of tasks and requirements.

\begin{lstlisting}[caption={step11\_evaluation\_suite.py}, label={lst:step11}]
"""
PHASE 4 - STEP 11: LLM-as-Judge Evaluation Suite
Harness Pillar: Evaluation & Guardrails
Goal: Build a pipeline that tests your agent and grades results.
"""

from pydantic import BaseModel
from typing import Literal

class EvalResult(BaseModel):
    verdict: Literal["pass", "fail"]
    score: float
    reasoning: str

def llm_judge(task: str, output: str, requirements: list) -> EvalResult:
    response = client.messages.parse(
        model="claude-opus-4-5-20251101",
        output_format=EvalResult,
        system="You are a strict technical evaluator.",
        messages=[{"role": "user", "content": f"Task: {task}\nOutput: {output}\nReqs: {requirements}"}],
    )
    return response.content[0].parsed
\end{lstlisting}

\subsection*{Key Takeaways}
\begin{itemize}
    \item The evaluator model should be separate from the agent model to maintain objectivity.
    \item Use structured outputs (Pydantic) for the judge to ensure deterministic verdicts and parseable scores.
    \item Run evaluation suites in CI/CD before deploying any agent updates.
\end{itemize}

\section{Step 12: Human-in-the-Loop (HITL)}
For high-stakes actions like sending emails or mutating databases, a human approval gate is mandatory. The harness pauses execution and waits for a human decision (Approve, Reject, or Modify).

\begin{lstlisting}[caption={step12\_hitl\_agent.py}, label={lst:step12}]
"""
PHASE 4 - STEP 12: Human-in-the-Loop (HITL) Approval Gate
Harness Pillar: Evaluation & Guardrails
Goal: Pause the agent on high-stakes actions.
"""

HIGH_STAKES_TOOLS = {"send_email", "write_database"}

def run_hitl_agent(task: str):
    messages = [{"role": "user", "content": task}]
    # ... ReAct Loop ...
    if tool_name in HIGH_STAKES_TOOLS:
        print(f"[APPROVAL REQUIRED]: {tool_name}")
        decision = "approve" # Simulate human input
        if decision == "reject":
            # Return error to model
            pass
\end{lstlisting}

\subsection*{Key Takeaways}
\begin{itemize}
    \item Maintain an explicit allowlist of tools that require a permission bridge.
    \item The approval gate pauses the agentic loop rather than canceling it, allowing the human to course-correct the input.
    \item Real-world implementations often use Slack buttons, email approval links, or specialized web modals.
\end{itemize}
